% =====================================================================
% S214-Core: 3-adic Accessibility Potentials
% Standalone technical document (companion to s216_paper.tex).
% Compile on Overleaf or with: pdflatex s214_core.tex
% =====================================================================
\documentclass[11pt,a4paper]{article}

\usepackage[T1]{fontenc}
\usepackage[utf8]{inputenc}
\usepackage{lmodern}
\usepackage[english]{babel}

\usepackage{amsmath,amssymb,amsthm}
\usepackage{mathtools}
\usepackage{microtype}
\usepackage[margin=2.5cm]{geometry}
\usepackage{hyperref}
\hypersetup{
    colorlinks=true,
    linkcolor=blue!50!black,
    citecolor=blue!50!black,
    urlcolor=blue!50!black
}

% --- Theorem environments -------------------------------------------------
\theoremstyle{plain}
\newtheorem{theorem}{Theorem}[section]
\newtheorem{lemma}[theorem]{Lemma}
\newtheorem{proposition}[theorem]{Proposition}
\newtheorem{corollary}[theorem]{Corollary}

\theoremstyle{definition}
\newtheorem{definition}[theorem]{Definition}

\theoremstyle{remark}
\newtheorem{conjecture}[theorem]{Conjecture}
\newtheorem{remark}[theorem]{Remark}

% --- Title ---------------------------------------------------------------
\title{S214-Core: 3-adic Accessibility Potentials\\
\large Companion technical note to the S216 closure-to-barrier theorem}
\author{Anonymous (double-blind submission)}
\date{}

\begin{document}
\maketitle

\begin{abstract}
We isolate the structural mechanism behind the accessibility barrier
associated with the S202 subcritical block of the admissible inverse
Collatz tree. We introduce the root and negative-run \(3\)-adic
proximities \(\rho_+,\rho_-\) and their associated deficit
\(\delta_+\), prove an elementary separation lemma stating that a
\(3\)-adic cylinder cannot be simultaneously close to \(1\) and to
\(-1\), and derive an arithmetic constraint forcing negative defect
runs to live near \(-1\). We then formalise candidate accessibility
potentials \(\Phi(j,c)\) that bound the cost
\(\mathcal C^{\leftarrow}_{S202}(m,Q)\) via Bellman--Ford-style edge
inequalities, state the S202 accessibility barrier conjecture and its
universal generalisation, and connect the framework to a long-term
reduction strategy for the Collatz problem. All lemmas in this note are
proved unconditionally; the conjectures are stated explicitly as open
problems.
\end{abstract}

\tableofcontents

% =====================================================================
% S214-Core: 3-adic accessibility potentials
% Standalone section fragment. \input from a host document that defines
% the lemma/proposition/conjecture/proof environments and loads amsmath.
% =====================================================================

\section{S214-Core: 3-adic Accessibility Potentials}
\label{sec:S214-core}

The purpose of this section is to isolate the structural mechanism behind
the accessibility barrier for the S202 block. The block S202 shows that
local subcriticality in the admissible inverse system is possible:
\[
A(w_{202})-2q(w_{202})=-1.
\]
Thus the naive inequality \(A\ge 2q\) is false. The correct question is no
longer whether subcritical blocks exist, but whether such blocks are
accessible from the root with enough gain to overcome their entrance cost.

The guiding principle is:

\[
\boxed{
\text{Subcritical }3\text{-adic blocks may exist, but they may be protected
by accessibility barriers.}
}
\]

The first test case is the S202 block.

% ---------------------------------------------------------------------
\subsection{The S202 fixed cylinder}
\label{subsec:S202-cylinder}

Let
\[
w_{202}=10100000001000010000000000,
\]
with certified data
\[
A_{202}=43,\qquad q_{202}=22,\qquad D_{202}=A_{202}-2q_{202}=-1.
\]
The associated affine action is
\[
F_{202}(x)=\frac{2^{43}x-B_{202}}{3^{22}},
\]
where
\[
B_{202}=919447060349.
\]
Its formal fixed point in \(\mathbb Z_3\) is
\[
a_{202}
=
\frac{B_{202}}{2^{43}-3^{22}}.
\]

A direct computation gives
\[
a_{202}\equiv 1+3^{22}\pmod{3^{24}}.
\]
Therefore
\[
v_3(a_{202}-1)=22.
\]

For \(m\ge1\), define
\[
R_m=22m+2
\]
and the S202 target cylinder
\[
C_m
=
a_{202}+3^{R_m}\mathbb Z_3.
\]

The accessibility cost is
\[
\mathcal C^{\leftarrow}_{S202}(m,Q)
=
\inf
\left\{
D(u):
q(u)\le Q,\
n(u)\equiv a_{202}\pmod{3^{R_m}}
\right\}.
\]

Since concatenating \(m\) copies of \(w_{202}\) decreases the defect by \(m\),
we have
\[
D(u w_{202}^{m})
=
D(u)-m.
\]
Hence, if
\[
\mathcal C^{\leftarrow}_{S202}(m,Q)<m,
\]
then there exists an accessible concatenation with negative defect:
\[
D(u w_{202}^{m})<0.
\]

Conversely, if
\[
\mathcal C^{\leftarrow}_{S202}(m,Q)\ge m,
\]
then S202 cannot produce a net subcritical concatenation under the budget
\(q(u)\le Q\).

% ---------------------------------------------------------------------
\subsection{Root and negative-run proximities}
\label{subsec:rho-plus-minus}

Let \([c]_r\) denote the \(3\)-adic cylinder
\[
[c]_r=c+3^r\mathbb Z_3.
\]

We define the root proximity
\[
\rho_+(c,r)
=
\min\{v_3(c-1),r\},
\]
and the root deficit
\[
\delta_+(c,r)
=
r-\rho_+(c,r).
\]

Thus
\[
\delta_+(c,r)=0
\]
if and only if
\[
[c]_r
\]
contains the root \(1\).

We also define the negative-run proximity
\[
\rho_-(c,r)
=
\min\{v_3(c+1),r\}.
\]

The quantity \(\rho_-(c,r)\) measures proximity to the \(3\)-adic point
\(-1\), which controls the possibility of long negative defect runs.

% ---------------------------------------------------------------------
\subsection{Separation between the root world and the negative-run world}
\label{subsec:root-negative-separation}

\begin{lemma}[Root/negative separation]
\label{lem:root-negative-separation}
For every integer \(c\) and every \(r\ge1\),
\[
\min\bigl(\rho_+(c,r),\rho_-(c,r)\bigr)=0.
\]
Equivalently, a \(3\)-adic cylinder cannot be simultaneously close to
\(1\) and to \(-1\).
\end{lemma}

\begin{proof}
Suppose first that \(\rho_+(c,r)>0\). Then
\[
v_3(c-1)>0,
\]
so
\[
c\equiv1\pmod3.
\]
Therefore
\[
c+1\equiv2\pmod3,
\]
and hence
\[
v_3(c+1)=0.
\]
Thus
\[
\rho_-(c,r)=0.
\]

Similarly, if \(\rho_-(c,r)>0\), then
\[
c\equiv -1\pmod3,
\]
so
\[
c-1\equiv -2\not\equiv0\pmod3.
\]
Hence
\[
v_3(c-1)=0,
\]
and therefore
\[
\rho_+(c,r)=0.
\]

Thus the two truncated valuations cannot both be positive.
\end{proof}

This elementary separation is one of the central structural facts of the
program. It means that paths attempting to exploit negative defect gains
must move into the \(-1\)-world, whereas paths attempting to terminate at
the root must move into the \(1\)-world. These two regions are disjoint
already modulo \(3\).

% ---------------------------------------------------------------------
\subsection{Negative defect steps and proximity to \(-1\)}
\label{subsec:negative-runs}

Recall that the defect is
\[
D=A-2q.
\]
Under a branch \(1\),
\[
D\mapsto D+\tau(n),
\]
while under a branch \(0\),
\[
D\mapsto D+\tau(n)-2.
\]

The only negative defect steps occur when
\[
\tau(n)=1
\]
and the branch is \(0\). This happens precisely for
\[
n\equiv2,8\pmod9.
\]

In that case the map is
\[
L(n)=\frac{2n-1}{3}.
\]

\begin{lemma}[Iterated negative branch]
\label{lem:iterated-negative-branch}
For every \(k\ge0\),
\[
L^k(n)
=
\frac{2^k n+2^k-3^k}{3^k}.
\]
Equivalently,
\[
L^k(n)+1
=
\frac{2^k(n+1)}{3^k}.
\]
\end{lemma}

\begin{proof}
For \(k=0\), the formula is immediate. Assume it holds for \(k\). Then
\[
L^{k+1}(n)
=
L(L^k(n))
=
\frac{2L^k(n)-1}{3}.
\]
Using the induction hypothesis,
\[
2L^k(n)-1
=
2\cdot
\frac{2^k n+2^k-3^k}{3^k}
-1.
\]
Thus
\[
2L^k(n)-1
=
\frac{2^{k+1}n+2^{k+1}-2\cdot3^k-3^k}{3^k}
=
\frac{2^{k+1}n+2^{k+1}-3^{k+1}}{3^k}.
\]
Dividing by \(3\), we obtain
\[
L^{k+1}(n)
=
\frac{2^{k+1}n+2^{k+1}-3^{k+1}}{3^{k+1}}.
\]
The equivalent formula follows by adding \(1\):
\[
L^k(n)+1
=
\frac{2^k n+2^k}{3^k}
=
\frac{2^k(n+1)}{3^k}.
\]
\end{proof}

\begin{lemma}[Negative runs force proximity to \(-1\)]
\label{lem:negative-runs-force-minus-one}
If \(k\) consecutive negative steps are defined starting from \(n\), then
\[
n\equiv -1\pmod{3^k}.
\]
Equivalently,
\[
v_3(n+1)\ge k.
\]
\end{lemma}

\begin{proof}
If \(k\) consecutive applications of
\[
L(n)=\frac{2n-1}{3}
\]
are defined over the integers, then \(L^k(n)\in\mathbb Z\). By
Lemma~\ref{lem:iterated-negative-branch},
\[
L^k(n)+1
=
\frac{2^k(n+1)}{3^k}.
\]
Since \(2^k\) is invertible modulo powers of \(3\), integrality forces
\[
3^k\mid n+1.
\]
Hence
\[
n\equiv -1\pmod{3^k}.
\]
\end{proof}

The converse requires the fine admissibility condition: integrality alone
is not sufficient to guarantee that all intermediate states remain in the
negative classes \(2,8\pmod9\). More precisely, a genuine negative run of
length \(k\) corresponds to two of the three classes above
\[
-1\pmod{3^k}.
\]
Equivalently, one needs
\[
n\equiv -1+3^k m\pmod{3^{k+1}},
\]
with
\[
m\in\{0,1,2\},
\qquad
m\not\equiv 2^k\pmod3.
\]

This refinement is important computationally, but the coarse implication
\[
\text{negative run of length }k
\Longrightarrow
v_3(n+1)\ge k
\]
is already enough to justify the appearance of \(\rho_-\) in the potential.

% ---------------------------------------------------------------------
\subsection{The S202 root-distance deficit}
\label{subsec:S202-root-distance-deficit}

\begin{lemma}[S202 root-distance deficit]
\label{lem:S202-root-deficit}
For the S202 fixed point and
\[
R_m=22m+2,
\]
one has
\[
\rho_+(a_{202},R_m)=22,
\]
and therefore
\[
\delta_+(a_{202},R_m)
=
R_m-22
=
22m-20.
\]
\end{lemma}

\begin{proof}
The certified computation gives
\[
a_{202}\equiv1+3^{22}\pmod{3^{24}}.
\]
Thus
\[
a_{202}-1
=
3^{22}u
\]
for some \(3\)-adic unit \(u\), since the coefficient of \(3^{22}\) is
nonzero modulo \(3\). Hence
\[
v_3(a_{202}-1)=22.
\]
For \(m\ge1\),
\[
R_m=22m+2\ge24
\]
except for the minimal boundary where the same valuation conclusion remains
valid by truncation. Therefore
\[
\rho_+(a_{202},R_m)
=
\min\{22,R_m\}
=
22.
\]
Consequently,
\[
\delta_+(a_{202},R_m)
=
R_m-\rho_+(a_{202},R_m)
=
22m+2-22
=
22m-20.
\]
\end{proof}

This lemma is the first quantitative indication of an accessibility
barrier. Although \(a_{202}\) is close to the root \(1\) modulo \(3^{22}\),
the cylinder \(C_m\) requires agreement modulo
\[
3^{22m+2}.
\]
Thus the missing root precision is
\[
22m-20.
\]

The S202 block gives only \(m\) units of defect gain after \(m\)
concatenations. Hence the accessibility problem asks whether the system can
correct approximately \(22m\) missing \(3\)-adic digits while paying less
than \(m\) units of defect.

% ---------------------------------------------------------------------
\subsection{Candidate accessibility potentials}
\label{subsec:candidate-potentials}

The previous lemmas suggest that any useful path toward the S202 cylinder
must balance two competing mechanisms:

\[
\boxed{
\text{approaching the root }1
}
\]
and
\[
\boxed{
\text{exploiting negative runs near }-1.
}
\]

This motivates potentials of the form
\[
\Phi(j,c)
=
\lambda\,\delta_+(c,r)
+
\gamma\,\rho_-(c,r)
-
\mu(Q-j)
+
\psi(c\bmod 3^h),
\]
where
\[
r=R_m+j.
\]

Here:

\begin{itemize}
    \item \(j\) is the number of inverse zero-steps already used;
    \item \(Q-j\) is the remaining zero budget;
    \item \(\delta_+(c,r)\) measures distance from the root \(1\);
    \item \(\rho_-(c,r)\) measures proximity to the negative-run region
          around \(-1\);
    \item \(\psi(c\bmod3^h)\) is a finite local correction.
\end{itemize}

The role of the term
\[
-\mu(Q-j)
\]
is to account for the maximum possible future defect decrease, since each
remaining zero can reduce the defect by at most \(1\).

The role of
\[
\gamma\rho_-(c,r)
\]
is to penalize excessive reliance on negative runs: long negative runs
require proximity to \(-1\), but by
Lemma~\ref{lem:root-negative-separation}, this is incompatible with
proximity to the root \(1\).

The finite correction
\[
\psi(c\bmod3^h)
\]
is necessary because the evolution of \(\rho_+\) and \(\rho_-\) under the
inverse branches is not determined solely by their current values. Local
residue data modulo \(3^h\) is required.

% ---------------------------------------------------------------------
\subsection{Dual certificate condition}
\label{subsec:dual-certificate-condition}

Let \(G_{m,Q}\) be the inverse cylinder graph associated with the S202
target cylinder and zero budget \(Q\). A vertex is written
\[
v=(j,c),
\]
and each edge
\[
v\to v'
\]
has weight
\[
\omega(v,v')\in\{-1,0,1,2,3,4\}.
\]

A potential \(\Phi\) certifies a lower bound if:

\[
\Phi(g)\le0
\]
for every goal vertex \(g\), and
\[
\Phi(v)\le \omega(v,v')+\Phi(v')
\]
for every edge \(v\to v'\).

\begin{proposition}[Dual barrier certificate]
\label{prop:dual-barrier-certificate}
If a potential \(\Phi\) satisfies
\[
\Phi(g)\le0
\]
for every goal vertex \(g\), and
\[
\Phi(v)\le \omega(v,v')+\Phi(v')
\]
for every edge \(v\to v'\), then every path from \(s\) to a goal has weight
at least
\[
\Phi(s).
\]
In particular, if
\[
\Phi(s)\ge T,
\]
then
\[
\mathcal C^{\leftarrow}_{S202}(m,Q)\ge T.
\]
\end{proposition}

\begin{proof}
Let
\[
s=v_0\to v_1\to\cdots\to v_k=g
\]
be a path from the start vertex \(s\) to a goal vertex \(g\). Applying the
edge inequality to every edge gives
\[
\Phi(v_i)
\le
\omega(v_i,v_{i+1})+\Phi(v_{i+1}).
\]
Summing telescopically,
\[
\Phi(s)
\le
\sum_{i=0}^{k-1}\omega(v_i,v_{i+1})+\Phi(g).
\]
Since \(g\) is a goal and \(\Phi(g)\le0\), we get
\[
\Phi(s)
\le
\sum_{i=0}^{k-1}\omega(v_i,v_{i+1}).
\]
Thus every path from \(s\) to a goal has weight at least \(\Phi(s)\).
\end{proof}

This proposition is the abstract form of the S216 closure-to-barrier
certificate. The Bellman--Ford closure certificate is a finite,
machine-checkable way of constructing such a potential or, equivalently,
proving that no path of weight \(<T\) exists.

% ---------------------------------------------------------------------
\subsection{S202 accessibility barrier conjecture}
\label{subsec:S202-barrier-conjecture}

The computations for S202 suggest the following conjecture.

\begin{conjecture}[S202 accessibility barrier]
\label{conj:S202-accessibility-barrier}
For the S202 block, there exists a constant \(K>0\) such that, for all
sufficiently large \(m\),
\[
\mathcal C^{\leftarrow}_{S202}(m,Km)\ge m.
\]
Equivalently, no access path using at most \(Km\) inverse zero-steps can
enter the cylinder
\[
a_{202}+3^{22m+2}\mathbb Z_3
\]
with defect smaller than the \(m\) units of defect gained by
\(w_{202}^{m}\).
\end{conjecture}

A stronger version would assert that there exists \(\varepsilon>0\) such
that
\[
\mathcal C^{\leftarrow}_{S202}(m,Km)
\ge
(1+\varepsilon)m
\]
for all sufficiently large \(m\).

The meaning of the conjecture is:

\[
\boxed{
\text{S202 is a subcritical }3\text{-adic block, but it is not cheaply
accessible from the root.}
}
\]

% ---------------------------------------------------------------------
\subsection{Universal subcritical accessibility barrier}
\label{subsec:universal-subcritical-barrier}

The S202 block should be regarded as the first test case of a broader
principle.

Let \(w\) be any admissible block with data
\[
A(w),\qquad q(w),\qquad B(w),
\]
and defect
\[
D(w)=A(w)-2q(w).
\]

Suppose
\[
D(w)<0.
\]
Write
\[
\Delta(w)=-D(w)>0.
\]

If \(w\) has a \(3\)-adic fixed point
\[
a_w=
\frac{B(w)}{2^{A(w)}-3^{q(w)}}\in\mathbb Z_3,
\]
then the natural analogue of the S202 barrier is
\[
\mathcal C^{\leftarrow}_{w}(m,Q)
\ge
m\Delta(w).
\]

\begin{conjecture}[Universal subcritical accessibility barrier]
\label{conj:universal-subcritical-barrier}
For every admissible subcritical block \(w\) with \(D(w)<0\), there exists a
linear budget \(Q_m=O(m)\) such that, for all sufficiently large \(m\),
\[
\mathcal C^{\leftarrow}_{w}(m,Q_m)
\ge
-mD(w).
\]
Equivalently, the cost of accessing the \(m\)-th fixed cylinder of a
subcritical block is at least the total defect gain obtained by repeating
that block \(m\) times.
\end{conjecture}

This conjecture expresses the general principle:

\[
\boxed{
\text{local subcriticality does not imply globally accessible
subcriticality.}
}
\]

% ---------------------------------------------------------------------
\subsection{Relation to the Collatz problem}
\label{subsec:relation-to-collatz}

The S202 accessibility barrier, even if proved, does not by itself prove
the Collatz conjecture. It only proves that one specific subcritical
\(3\)-adic block cannot be exploited cheaply from the root.

To obtain a route toward Collatz, one needs an additional reduction
principle.

\begin{conjecture}[Subcritical reduction of counterexamples]
\label{conj:subcritical-reduction}
If the Collatz conjecture fails, then there exists an accessible family of
admissible inverse words whose asymptotic defect becomes negative after
entering a recurrent subcritical block or family of blocks.
\end{conjecture}

In schematic form, the desired strategy is:

\[
\boxed{
\text{Counterexample}
\Longrightarrow
\text{accessible subcritical family}
}
\]
and
\[
\boxed{
\text{universal accessibility barrier}
\Longrightarrow
\text{no accessible subcritical family}.
}
\]

Together these would imply the nonexistence of a counterexample.

Thus the long-term program is:

\[
\boxed{
\text{Reduction of counterexamples}
+
\text{universal subcritical accessibility barrier}
\Longrightarrow
\text{Collatz}.
}
\]

At the current stage, S202 provides the first nontrivial test case for this
program, and S214 provides the potential-theoretic language needed to move
from finite certificates to structural barriers.

% ---------------------------------------------------------------------
\subsection{Immediate formal milestones}
\label{subsec:S214-formal-milestones}

The immediate formal targets are:

\begin{enumerate}
    \item Complete the S216 closure-to-barrier bridge in Lean.
    \item Obtain the first formal theorem:
    \[
    \mathcal C^{\leftarrow}_{S202}(1,3)\ge1.
    \]
    \item Scale the certificate generation to:
    \[
    (m,Q,T)=(1,5,1),\ (1,8,1),\ (1,10,1),\ (2,6,2).
    \]
    \item Extract empirical profiles of the cheap region using:
    \[
    \rho_+,\qquad \rho_-,\qquad j,\qquad D.
    \]
    \item Use these profiles to search for explicit finite-depth potentials
    of the form
    \[
    \Phi(j,c)
    =
    \lambda\delta_+(c,r)
    +
    \gamma\rho_-(c,r)
    -
    \mu(Q-j)
    +
    \psi(c\bmod3^h).
    \]
\end{enumerate}

The first fully verified Lean theorem
\[
\mathcal C^{\leftarrow}_{S202}(1,3)\ge1
\]
will constitute the first formal accessibility barrier in the S202 chain.


\end{document}
